\begin{tomtat}
	\begin{minipage}[t]{0.3\textwidth}
		\vspace{0pt}
		\begin{tomtat}
			\begin{center}
				\textbf{ĐỀ 51}
			\end{center}
			
		\end{tomtat}
	\end{minipage}
	\hfill
	\begin{minipage}[t]{0.5\textwidth}
		\vspace{0pt}
		\begin{tomtat}
			\begin{center}
				\includegraphics[scale=0.08]{images/51}
			\end{center}
		\end{tomtat}
	\end{minipage}
\end{tomtat}
\textbf{I. PHẦN TRẮC NGHIỆM}
\begin{ex} % Câu 1
	Giá trị của biểu thức $A=\sqrt{18}+\sqrt[3]{64}-\sqrt{50}$ bằng:
	\choice{$4-2\sqrt{2}$}{$-4-2\sqrt{2}$}{$-4+2\sqrt{2}$}{$4+2\sqrt{2}$}
	\loigiai{
		\textbf{Đáp án đúng là A.} \\
		\begin{align*}
			A &= \sqrt{18}+\sqrt[3]{64}-\sqrt{50} \\
			&= \sqrt{9 \cdot 2} + \sqrt[3]{4^3} - \sqrt{25 \cdot 2} \\
			&= 3\sqrt{2} + 4 - 5\sqrt{2} \\
			&= 4 + (3\sqrt{2} - 5\sqrt{2}) \\
			&= 4 - 2\sqrt{2}
		\end{align*}
	}
\end{ex}

\begin{ex} % Câu 2
	Hình trụ có bán kính đáy bằng $2\sqrt{3}$cm và thể tích bằng $24\pi \text{ cm}^3$. Chiều cao của hình trụ này là:
	\choice{$h=2 \text{ cm}$}{$h=6 \text{ cm}$}{$h=2\sqrt{3}\text{cm}$}{$h=1,5\text{cm}$}
	\loigiai{
		\textbf{Đáp án đúng là A.} \\
		Bán kính đáy $R = 2\sqrt{3} \text{ cm}$.
		Thể tích hình trụ $V = 24\pi \text{ cm}^3$.
		Công thức tính thể tích hình trụ là $V = \pi R^2 h$.
		Ta có:
		\begin{align*}
			24\pi &= \pi (2\sqrt{3})^2 h \\
			24\pi &= \pi (4 \cdot 3) h \\
			24\pi &= 12\pi h \\
			h &= \frac{24\pi}{12\pi} \\
			h &= 2 \text{ (cm)}
		\end{align*}
	}
\end{ex}

\begin{ex} % Câu 3
	Bất phương trình $3x-5>4x+2$ có nghiệm là:
	\choice{$x>7$}{$x<7$}{$x<-7$}{$x>-7$}
	\loigiai{
		\textbf{Đáp án đúng là C.} \\
		\begin{align*}
			3x-5 &> 4x+2 \\
			3x-4x &> 2+5 \\
			-x &> 7 \\
			x &< -7 \quad \text{(Khi nhân/chia cho số âm, đổi chiều bất phương trình)}
		\end{align*}
	}
\end{ex}

\begin{ex} % Câu 4
	Đồ thị hàm số $y=-\frac{1}{\sqrt{3}}x^2$ đi qua điểm nào sau đây?
	\choice{$(-1;\sqrt{3})$}{$(-\sqrt{3};-\sqrt{3})$}{$(3;\sqrt{3})$}{$(-1;3)$}
	\loigiai{
		\textbf{Đáp án đúng là B.} \\
		Để kiểm tra một điểm có thuộc đồ thị hàm số hay không, ta thay tọa độ của điểm đó vào phương trình hàm số $y=-\frac{1}{\sqrt{3}}x^2$.
		A. Với $(-1;\sqrt{3})$: Thay $x=-1, y=\sqrt{3} \Rightarrow \sqrt{3} = -\frac{1}{\sqrt{3}}(-1)^2 \Rightarrow \sqrt{3} = -\frac{1}{\sqrt{3}}$ (Sai).
		B. Với $(-\sqrt{3};-\sqrt{3})$: Thay $x=-\sqrt{3}, y=-\sqrt{3} \Rightarrow -\sqrt{3} = -\frac{1}{\sqrt{3}}(-\sqrt{3})^2 \Rightarrow -\sqrt{3} = -\frac{1}{\sqrt{3}}(3) \Rightarrow -\sqrt{3} = -\frac{3}{\sqrt{3}} \Rightarrow -\sqrt{3} = -\sqrt{3}$ (Đúng).
		C. Với $(3;\sqrt{3})$: Thay $x=3, y=\sqrt{3} \Rightarrow \sqrt{3} = -\frac{1}{\sqrt{3}}(3)^2 \Rightarrow \sqrt{3} = -\frac{9}{\sqrt{3}} \Rightarrow \sqrt{3} = -3\sqrt{3}$ (Sai).
		D. Với $(-1;3)$: Thay $x=-1, y=3 \Rightarrow 3 = -\frac{1}{\sqrt{3}}(-1)^2 \Rightarrow 3 = -\frac{1}{\sqrt{3}}$ (Sai).
		Vậy, điểm $(-\sqrt{3};-\sqrt{3})$ thuộc đồ thị hàm số.
	}
\end{ex}

\begin{ex} % Câu 5
	Đo chiều cao (đơn vị: centimét) của học sinh lớp 9A ở một trường THCS người ta thu được bảng tần số ghép nhóm như sau:
	\begin{center}
		\begin{tabular}{|c|c|c|c|c|}
			\hline
			[150;158) & [158;161) & [161;164) & [164;167) \\
			\hline
			5 & 12 & 15 & 8 \\
			\hline
		\end{tabular}
	\end{center}
	Tỉ lệ học sinh có chiều cao từ 158 cm đến dưới 161 cm là:
	\choice{$12,5\%$}{$37,5\%$}{$20\%$}{$30\%$}
	\loigiai{
		\textbf{Đáp án đúng là D.} \\
		Tổng số học sinh được khảo sát là: $5 + 12 + 15 + 8 = 40$ học sinh.
		Số học sinh có chiều cao từ 158 cm đến dưới 161 cm (nhóm [158;161)) là 12 học sinh.
		Tỉ lệ học sinh có chiều cao từ 158 cm đến dưới 161 cm là:
		\begin{align*}
			\text{Tỉ lệ} &= \frac{\text{Số học sinh trong nhóm [158;161)}}{\text{Tổng số học sinh}} \times 100\% \\
			&= \frac{12}{40} \times 100\% \\
			&= \frac{3}{10} \times 100\% \\
			&= 0,3 \times 100\% \\
			&= 30\%
		\end{align*}
	}
\end{ex}
\begin{ex} % Câu 6
	Tổng các nghiệm của phương trình $(2x+1)(x-2)=0$ là:
	\choice{$\frac{5}{2}$}{$\frac{3}{2}$}{$-\frac{3}{2}$}{$-\frac{5}{2}$}
	\loigiai{
		\textbf{Đáp án đúng là B.} \\
		Phương trình tích $(2x+1)(x-2)=0$ xảy ra khi một trong hai thừa số bằng 0:
		\begin{itemize}
			\item $2x+1=0 \Rightarrow 2x = -1 \Rightarrow x_1 = -\frac{1}{2}$
			\item $x-2=0 \Rightarrow x_2 = 2$
		\end{itemize}
		Tổng các nghiệm của phương trình là:
		$x_1 + x_2 = -\frac{1}{2} + 2 = -\frac{1}{2} + \frac{4}{2} = \frac{3}{2}$.
	}
\end{ex}

\begin{ex} % Câu 7
	Cho $\triangle ABC$ vuông tại C có AB=6 cm, BC=4 cm. Tính $\sin A$?
	\choice{$\frac{\sqrt{5}}{3}$}{$\frac{6}{4}$}{$\frac{2}{3}$}{$\frac{3}{2}$}
	\loigiai{
		\textbf{Đáp án đúng là C.} \\
		Trong tam giác ABC vuông tại C:
		$\sin A = \frac{\text{cạnh đối}}{\text{cạnh huyền}} = \frac{BC}{AB}$
		\begin{align*}
			\sin A &= \frac{4}{6} \\
			&= \frac{2}{3}
		\end{align*}
	}
\end{ex}

\begin{ex} % Câu 8
	Một hộp đựng 36 tấm thẻ giống nhau được đánh số $1; 2; 3; 4; \dots; 35; 36$. Bạn Nam rút ngẫu nhiên một tấm thẻ trong hộp. Xác suất của biến cố rút được tấm thẻ ghi số nguyên tố là:
	\choice{$\frac{13}{36}$}{$\frac{1}{3}$}{$\frac{1}{4}$}{$\frac{11}{36}$}
	\loigiai{
		\textbf{Đáp án đúng là D.} \\
		Tổng số tấm thẻ trong hộp là 36. Vậy, số phần tử của không gian mẫu là $n(\Omega) = 36$.
		Gọi A là biến cố "rút được tấm thẻ ghi số nguyên tố".
		Các số nguyên tố từ 1 đến 36 là:
		$2, 3, 5, 7, 11, 13, 17, 19, 23, 29, 31$.
		Có 11 số nguyên tố trong khoảng này.
		Vậy, số phần tử của biến cố A là $n(A) = 11$.
		Xác suất của biến cố A là:
		$P(A) = \frac{n(A)}{n(\Omega)} = \frac{11}{36}$.
	}
\end{ex}

\newpage
\textbf{II. PHẦN TỰ LUẬN}
\begin{ex} % Câu 9
	Rút gọn biểu thức: $P=\left(\frac{\sqrt{x}+2}{x-4} + \frac{2}{x-2\sqrt{x}}\right) : \frac{\sqrt{x}+2}{\sqrt{x}-2}$ (với $x>0, x \ne 4$)
	\loigiai{
		Điều kiện: $x>0, x \ne 4$.
		Ta có:
		$x-4 = (\sqrt{x}-2)(\sqrt{x}+2)$
		$x-2\sqrt{x} = \sqrt{x}(\sqrt{x}-2)$
		Xét phần trong ngoặc:
		\begin{align*}
			\frac{\sqrt{x}+2}{x-4} + \frac{2}{x-2\sqrt{x}} &= \frac{\sqrt{x}+2}{(\sqrt{x}-2)(\sqrt{x}+2)} + \frac{2}{\sqrt{x}(\sqrt{x}-2)} \\
			&= \frac{1}{\sqrt{x}-2} + \frac{2}{\sqrt{x}(\sqrt{x}-2)} \\
			&= \frac{\sqrt{x} + 2}{\sqrt{x}(\sqrt{x}-2)}
		\end{align*}
		Vậy biểu thức $P$ trở thành:
		\begin{align*}
			P &= \frac{\sqrt{x}+2}{\sqrt{x}(\sqrt{x}-2)} : \frac{\sqrt{x}+2}{\sqrt{x}-2} \\
			&= \frac{\sqrt{x}+2}{\sqrt{x}(\sqrt{x}-2)} \cdot \frac{\sqrt{x}-2}{\sqrt{x}+2}
		\end{align*}
		Vì $x>0$ nên $\sqrt{x}+2 \ne 0$. Vì $x \ne 4$ nên $\sqrt{x}-2 \ne 0$. Ta có thể rút gọn:
		\begin{align*}
			P &= \frac{1}{\sqrt{x}}
		\end{align*}
		Vậy, với $x>0, x \ne 4$ thì $P = \frac{1}{\sqrt{x}}$.
	}
\end{ex}

\begin{ex} % Câu 10
	Biết phương trình $2x^2-8x-5=0$ có hai nghiệm $x_1, x_2$. Không giải phương trình, tính giá trị của biểu thức $B=\frac{x_1}{x_2-1}+\frac{x_2}{x_1-1}+3$.
	\loigiai{
		Phương trình đã cho là $2x^2-8x-5=0$.
		Theo hệ thức Vi-ét, ta có:
		$ \begin{cases} x_1 + x_2 = -\frac{-8}{2} = 4 \\ x_1 x_2 = \frac{-5}{2} \end{cases} $
		Ta biến đổi biểu thức B:
		\begin{align*}
			B &= \frac{x_1(x_1-1) + x_2(x_2-1)}{(x_2-1)(x_1-1)} + 3 \\
			&= \frac{x_1^2-x_1 + x_2^2-x_2}{x_1x_2 - x_1 - x_2 + 1} + 3 \\
			&= \frac{(x_1^2+x_2^2) - (x_1+x_2)}{x_1x_2 - (x_1+x_2) + 1} + 3 \\
			&= \frac{((x_1+x_2)^2 - 2x_1x_2) - (x_1+x_2)}{x_1x_2 - (x_1+x_2) + 1} + 3
		\end{align*}
		Thay giá trị từ hệ thức Vi-ét vào biểu thức B:
		\begin{align*}
			B &= \frac{((4)^2 - 2(-\frac{5}{2})) - (4)}{(-\frac{5}{2}) - (4) + 1} + 3 \\
			&= \frac{(16 + 5) - 4}{-\frac{5}{2} - 4 + 1} + 3 \\
			&= \frac{21 - 4}{-\frac{5}{2} - 3} + 3 \\
			&= \frac{17}{-\frac{5}{2} - \frac{6}{2}} + 3 \\
			&= \frac{17}{-\frac{11}{2}} + 3 \\
			&= 17 \cdot \left(-\frac{2}{11}\right) + 3 \\
			&= -\frac{34}{11} + \frac{33}{11} \\
			&= -\frac{1}{11}
		\end{align*}
		Vậy, giá trị của biểu thức $B = -\frac{1}{11}$.
	}
\end{ex}

\begin{ex} % Câu 11
	Giải hệ phương trình: $ \begin{cases} 2x-5y=19 \\ 3x+y=3 \end{cases} $
	\loigiai{
		Ta có hệ phương trình:
		$ \begin{cases} 2x-5y=19 \quad &(1) \\ 3x+y=3 \quad &(2) \end{cases} $
		Từ phương trình (2), suy ra $y = 3-3x$.
		Thế $y = 3-3x$ vào phương trình (1):
		\begin{align*}
			2x - 5(3-3x) &= 19 \\
			2x - 15 + 15x &= 19 \\
			17x &= 19 + 15 \\
			17x &= 34 \\
			x &= 2
		\end{align*}
		Thay $x=2$ vào $y = 3-3x$:
		\begin{align*}
			y &= 3 - 3(2) \\
			y &= 3 - 6 \\
			y &= -3
		\end{align*}
		Vậy, nghiệm của hệ phương trình là $(x;y) = (2;-3)$.
	}
\end{ex}
\begin{ex} % Câu 12
	Một hộp có 20 viên bi với kích thước và khối lượng như nhau. Bạn An viết lên các viên bi đó các số $1, 2, 3, \dots, 20$; hai viên bi khác nhau thì viết hai số khác nhau. Xét phép thử “Lấy ngẫu nhiên một viên bi trong hộp”. Tính xác suất của biến cố “Số xuất hiện trên viên bi được lấy ra chia 5 dư 1”.
	\loigiai{
		Không gian mẫu $\Omega$ là tập hợp các số được viết trên 20 viên bi.
		Các phần tử của không gian mẫu được liệt kê là:
		$\Omega = \{1, 2, 3, 4, 5, 6, 7, 8, 9, 10, 11, 12, 13, 14, 15, 16, 17, 18, 19, 20\}$.
		Số phần tử của không gian mẫu là $n(\Omega) = 20$.
		
		Gọi A là biến cố: “Số xuất hiện trên viên bi được lấy ra chia cho 5 dư 1”.
		Để tìm các phần tử của biến cố A, ta liệt kê các số từ 1 đến 20 mà khi chia cho 5 có số dư là 1:
		\begin{itemize}
			\item $1 = 5 \cdot 0 + 1$ (1 chia cho 5 dư 1)
			\item $6 = 5 \cdot 1 + 1$ (6 chia cho 5 dư 1)
			\item $11 = 5 \cdot 2 + 1$ (11 chia cho 5 dư 1)
			\item $16 = 5 \cdot 3 + 1$ (16 chia cho 5 dư 1)
		\end{itemize}
		Các số tiếp theo sẽ lớn hơn 20 (ví dụ: $5 \cdot 4 + 1 = 21$).
		Vậy, các phần tử của biến cố A là $A = \{1, 6, 11, 16\}$.
		Số phần tử của biến cố A là $n(A) = 4$.
		
		Xác suất của biến cố A là:
		$P(A) = \frac{n(A)}{n(\Omega)} = \frac{4}{20} = \frac{1}{5}$.
	}
\end{ex}
\begin{ex} % Câu 12
	Quãng đường từ A đến B dài 90 km. Một người đi xe máy từ A đến B, khi đến B người đó nghỉ 30 phút rồi quay về A với vận tốc lớn hơn vận tốc lúc đi là 9 km/h. Thời gian kể từ lúc bắt đầu đi từ A đến khi trở về đến A là 5 giờ. Tính vận tốc của người đó khi đi từ A đến B?
	\loigiai{
		Gọi $x$ (km/h) là vận tốc của người đó khi đi từ A đến B.
		Điều kiện: $x > 0$.
		
		Thời gian người đó đi từ A đến B là: $t_{AB} = \frac{90}{x}$ (giờ).
		
		Vận tốc khi quay về A là: $x+9$ (km/h).
		Thời gian người đó đi từ B về A là: $t_{BA} = \frac{90}{x+9}$ (giờ).
		
		Thời gian nghỉ tại B là 30 phút = $\frac{1}{2}$ giờ = 0,5 giờ.
		
		Tổng thời gian kể từ lúc bắt đầu đi từ A đến khi trở về đến A là 5 giờ.
		Ta có phương trình:
		$t_{AB} + \text{thời gian nghỉ} + t_{BA} = 5$
		\begin{align*}
			\frac{90}{x} + 0,5 + \frac{90}{x+9} &= 5 \\
			\frac{90}{x} + \frac{90}{x+9} &= 5 - 0,5 \\
			\frac{90}{x} + \frac{90}{x+9} &= 4,5 \\
			\frac{90}{x} + \frac{90}{x+9} &= \frac{9}{2}
		\end{align*}
		Chia cả hai vế cho 9:
		\begin{align*}
			\frac{10}{x} + \frac{10}{x+9} &= \frac{1}{2}
		\end{align*}
		Nhân cả hai vế với $2x(x+9)$ để khử mẫu (với $x \ne 0, x \ne -9$):
		\begin{align*}
			10 \cdot 2(x+9) + 10 \cdot 2x &= x(x+9) \\
			20(x+9) + 20x &= x^2 + 9x \\
			20x + 180 + 20x &= x^2 + 9x \\
			40x + 180 &= x^2 + 9x \\
			x^2 + 9x - 40x - 180 &= 0 \\
			x^2 - 31x - 180 &= 0
		\end{align*}
		Giải phương trình bậc hai $x^2 - 31x - 180 = 0$:
		$\Delta = b^2 - 4ac = (-31)^2 - 4(1)(-180)$
		\begin{align*}
			\Delta &= 961 + 720 \\
			\Delta &= 1681
		\end{align*}
		$\sqrt{\Delta} = \sqrt{1681} = 41$.
		Phương trình có hai nghiệm phân biệt:
		$x_1 = \frac{-(-31) + 41}{2 \cdot 1} = \frac{31 + 41}{2} = \frac{72}{2} = 36$ (Thỏa mãn điều kiện $x>0$)
		$x_2 = \frac{-(-31) - 41}{2 \cdot 1} = \frac{31 - 41}{2} = \frac{-10}{2} = -5$ (Không thỏa mãn điều kiện $x>0$)
		Vậy, vận tốc của người đó khi đi từ A đến B là 36 km/h.
	}
\end{ex}
\begin{ex} % Câu 13
	Cho $\triangle ABC$ nhọn có $AB > AC$, nội tiếp đường tròn $(O)$, các đường cao $AD, CE$ cắt nhau tại $H$ ($D \in BC, E \in AB$). Qua $O$ kẻ đường thẳng song song với $AD$, cắt $BC, AB, AC$ lần lượt tại $M, N$ và $P$.
	\begin{enumerate}[a)]
		\item Chứng minh 4 điểm $A, C, D, E$ nằm trên một đường tròn.
		\item Gọi $I$ là trung điểm của $PN$. Chứng minh rằng $AP \cdot BM = PI \cdot HB$ và $IA \perp MH$.
	\end{enumerate}
	\loigiai{}
\end{ex}

\begin{ex} % Câu 14
	\begin{enumerate}[a)]
		\item Bác Tâm muốn xây một bể nước dạng hình hộp chữ nhật có nắp với dung tích 3000 lít. Đáy bể là một hình chữ nhật có chiều dài gấp đôi chiều rộng. Giá thuê công nhân để xây bể là 500 000 đồng cho mỗi mét vuông. Hỏi chi phí thấp nhất Bác Tâm cần bỏ ra để xây bể nước là bao nhiêu?
		\item Cho các số thực dương $a, b$ thỏa mãn điều kiện: $(a-b)\sqrt{ab} = \sqrt{2}$.
		Tìm giá trị nhỏ nhất của biểu thức: $P = \frac{(a^2+b^2)^2+1}{a^2+b^2}$.
	\end{enumerate}
	\loigiai{}
\end{ex}

\setcounterref{ex}{0}
